% GENERATED FILE. Edit canonical lessons, then run npm run generate:books.
% canonical-source-hash: fnv1a64:d7150a018c074cbd
% canonical-lessons: TE-C260-nunupaina, TE-C260-veccani, TE-C260-arudaina, TE-C260-vinta, TE-C260-dayagala

\chapter{Smooth, Warm, Rare, Strange, Kind}
\label{ch:te-smooth-warm-rare-strange-kind}
\hlchaptermodality{260}

\begin{chapteropening}
\textbf{By the end of this chapter:} \emph{I can say smooth, warm, rare, strange and kind.}

Complete the last lesson of chapter 260: I can say smooth, warm, rare, strange and kind.
\end{chapteropening}

\section[\texorpdfstring{\te{నునుపైన} (nunupaina)}{నునుపైన (nunupaina)}]{\te{నునుపైన} (nunupaina) — smooth}

\label{lesson:TE-C260-nunupaina}



\begin{quote}
Before the new one: say the Telugu for busy, then the Telugu for compulsory.
\end{quote}

\subsection*{You'll want to know: \te{నునుపైన}}
\textbf{\te{నునుపైన}} — \emph{nunupaina} — \textquotedblleft{}smooth\textquotedblright{}.

\begin{grammarlens}[title={What you've built}]
One more word for this chapter.
\end{grammarlens}

\subsection*{Your turn}
\begin{itemize}
\raggedright
  \item \emph{Say it:} \emph{nunupaina}
  \item \emph{Say it:} \emph{nunupaina}, once more
  \item \emph{Say it:} say \emph{nunupaina}
  \item \emph{Recall:} say the Telugu for busy, then the Telugu for compulsory, then say \emph{nunupaina} again
\end{itemize}

\begin{tcolorbox}[breakable,colback=teal!4,colframe=teal!35!black,title={Before you move on}]
What does \te{నునుపైన} mean? (\textquotedblleft{}Smooth\textquotedblright{}.) Say it once more.
\end{tcolorbox}
\section[\texorpdfstring{\te{వెచ్చని} (veccani)}{వెచ్చని (veccani)}]{\te{వెచ్చని} (veccani) — warm}

\label{lesson:TE-C260-veccani}



\begin{quote}
Before the new one: say the Telugu for compulsory, then the Telugu for smooth.
\end{quote}

\subsection*{You'll want to know: \te{వెచ్చని}}
\textbf{\te{వెచ్చని}} — \emph{veccani} — \textquotedblleft{}warm\textquotedblright{}.

\begin{grammarlens}[title={What you've built}]
One more word for this chapter.
\end{grammarlens}

\subsection*{Your turn}
\begin{itemize}
\raggedright
  \item \emph{Say it:} \emph{veccani}
  \item \emph{Say it:} \emph{veccani}, once more
  \item \emph{Say it:} say \emph{veccani}
  \item \emph{Recall:} say the Telugu for compulsory, then the Telugu for smooth, then say \emph{veccani} again
\end{itemize}

\begin{tcolorbox}[breakable,colback=teal!4,colframe=teal!35!black,title={Before you move on}]
What does \te{వెచ్చని} mean? (\textquotedblleft{}Warm\textquotedblright{}.) Say it once more.
\end{tcolorbox}
\section[\texorpdfstring{\te{అరుదైన} (arudaina)}{అరుదైన (arudaina)}]{\te{అరుదైన} (arudaina) — rare}

\label{lesson:TE-C260-arudaina}



\begin{quote}
Before the new one: say the Telugu for smooth, then the Telugu for warm.
\end{quote}

\subsection*{You'll want to know: \te{అరుదైన}}
\textbf{\te{అరుదైన}} — \emph{arudaina} — \textquotedblleft{}rare\textquotedblright{}.

\begin{grammarlens}[title={What you've built}]
One more word for this chapter.
\end{grammarlens}

\subsection*{Your turn}
\begin{itemize}
\raggedright
  \item \emph{Say it:} \emph{arudaina}
  \item \emph{Say it:} \emph{arudaina}, once more
  \item \emph{Say it:} say \emph{arudaina}
  \item \emph{Recall:} say the Telugu for smooth, then the Telugu for warm, then say \emph{arudaina} again
\end{itemize}

\begin{tcolorbox}[breakable,colback=teal!4,colframe=teal!35!black,title={Before you move on}]
What does \te{అరుదైన} mean? (\textquotedblleft{}Rare\textquotedblright{}.) Say it once more.
\end{tcolorbox}
\section[\texorpdfstring{\te{వింత} (vinta)}{వింత (vinta)}]{\te{వింత} (vinta) — strange, odd}

\label{lesson:TE-C260-vinta}



\begin{quote}
Before the new one: say the Telugu for warm, then the Telugu for rare.
\end{quote}

\subsection*{You'll want to know: \te{వింత}}
\textbf{\te{వింత}} — \emph{vinta} — \textquotedblleft{}strange, odd\textquotedblright{}.

\begin{grammarlens}[title={What you've built}]
One more word for this chapter.
\end{grammarlens}

\subsection*{Your turn}
\begin{itemize}
\raggedright
  \item \emph{Say it:} \emph{vinta}
  \item \emph{Say it:} \emph{vinta}, once more
  \item \emph{Say it:} say \emph{vinta}
  \item \emph{Recall:} say the Telugu for warm, then the Telugu for rare, then say \emph{vinta} again
\end{itemize}

\begin{tcolorbox}[breakable,colback=teal!4,colframe=teal!35!black,title={Before you move on}]
What does \te{వింత} mean? (\textquotedblleft{}Strange, odd\textquotedblright{}.) Say it once more.
\end{tcolorbox}
\section[\texorpdfstring{\te{దయగల} (dayagala)}{దయగల (dayagala)}]{\te{దయగల} (dayagala) — kind, compassionate}

\label{lesson:TE-C260-dayagala}



\begin{quote}
Before the new one: say the Telugu for rare, then the Telugu for strange.
\end{quote}

\subsection*{You'll want to know: \te{దయగల}}
\textbf{\te{దయగల}} — \emph{dayagala} — \textquotedblleft{}kind, compassionate\textquotedblright{}.

\begin{grammarlens}[title={What you've built}]
One more word for this chapter.
\end{grammarlens}

\subsection*{Your turn}
\begin{itemize}
\raggedright
  \item \emph{Say it:} \emph{dayagala}
  \item \emph{Say it:} \emph{dayagala}, once more
  \item \emph{Say it:} say \emph{dayagala}
  \item \emph{Recall:} say the Telugu for rare, then the Telugu for strange, then say \emph{dayagala} again
\end{itemize}

\begin{tcolorbox}[breakable,colback=teal!4,colframe=teal!35!black,title={Before you move on}]
What does \te{దయగల} mean? (\textquotedblleft{}Kind, compassionate\textquotedblright{}.) Say it once more.
\end{tcolorbox}
